This chapter presents the methodology and results for the sensitivity 
analysis, uncertainty quantification, and optimization of select 
transition scenarios explored in Chapters \ref{sec:fc_methods}-
\ref{sec:recycle_results}. This work is designed to identify how 
the variation of select hyperparameters affect select performance 
metrics of the fuel cycle and how the hyperparameters synergistically 
affect the performance metrics. This information is then used to design 
a transition scenario that is optimized for each of the performance metrics. 

\section{Methodology}
To perform the sensitivity analysis, uncertainty quantification, and 
transition optimization \Cyclus was coupled to Dakota \cite{adams_dakota_2019},
and Dakota was used as the driver for the analysis. The metrics of 
interest for this analysis include the amount of waste generated that 
must be sent to a repository, the mass of natural uranium required, and 
the amount of \gls{SWU} required. The first two performance 
metrics used in this work match metrics considered in the \gls{ES} 
\cite{wigeland_nuclear_2014}, and all three of these metrics were considered 
in previous sensitivity analysis when Dakota was coupled to \gls{DYMOND}
\cite{feng_sensitivity_2020,richards_application_2021}. 

\section{Sensitivity analysis results}

\section{Transition optimization results}